%!TEX root = ../dokumentation.tex

\chapter{Einordnung des Projekts in das Projektportfolio}

Die Einführung einer zentralen Lernplattform ist Teil der strategischen Weiterentwicklung
der Personalentwicklungs- und Weiterbildungsprozesse der dm-Gruppe und unterstützt
strategische Initiativen im Bereich Organisationsentwicklung, Digitalisierung und
Mitarbeiterqualifizierung. Die Portfolio-Steuerung stellt den strategischen Beitrag
und effizienten Ressourceneinsatz sicher. Die Lernplattform wurde als strategisch
relevantes Vorhaben priorisiert.

\section{Strategische Einordnung}

Die dm-Gruppe fördert die Entwicklung der Mitarbeitenden und digitalisiert sowie standardisiert interne Prozesse. Die Lernplattform stellt eine zentrale Infrastruktur für Weiterbildung, Schulungsmanagement und Kompetenzentwicklung bereit und verbessert durch Standardisierung von Buchung, Dokumentation und Reporting die Transparenz und Effizienz.

\begin{xltabular}{\textwidth}{@{}lX@{}}
    \caption{Strategischer Beitrag des Projekts} \label{tab:strategischer-beitrag} \\
    \hline
    \textbf{Strategische Zielsetzung}        & \textbf{Beitrag der Lernplattform}                          \\
    \hline
    Mitarbeiterentwicklung stärken           & Zentrale Plattform für Schulungen und Lerninhalte           \\

    Digitale Unterstützung interner Prozesse & Automatisierte Buchung, Dokumentation und Reporting         \\

    Transparenz über Qualifikationen         & Aggregierte Kompetenz- und Weiterbildungskennzahlen         \\

    Attraktivität als Arbeitgeber sichern    & Verbesserte und leichter zugängliche Weiterbildungsangebote \\
    \hline
\end{xltabular}

Das Projekt unterstützt damit den \ac{MOV} (Reduktion administrativer Aufwände, bessere Steuerbarkeit von Weiterbildungsmaßnahmen).

\section{Portfolio-Kontext}

Im Portfolio der dm-Gruppe werden Initiativen zur Umsetzung der Strategie gebündelt (u.\,a. Filialdigitalisierung, Supply Chain, datenbasierte Steuerung, \ac{HR}-Prozesse). Das Projekt „Lernplattform“ gehört zum \ac{HR}-Transformationscluster. Ziel ist die Modernisierung und konsistente Gestaltung von Personalentwicklung und Weiterbildung.

\begin{xltabular}{\textwidth}{@{}lXX@{}}
    \caption{Einordnung im Projektportfolio} \label{tab:portfolio-einordnung} \\
    \hline
    \textbf{Portfolio-Bereich} & \textbf{Beispielhafte Initiativen}        & \textbf{Strategischer Fokus}              \\
    \hline
    Filialdigitalisierung      & Self-Checkout, Filialanalyse              & Verbesserung des Kundenerlebnisses        \\

    Supply Chain               & Lagerautomatisierung, Dispositionssysteme & Effizienzsteigerung in Logistikprozessen  \\

    Datenplattform             & Data Warehouse, Analytics-Plattform       & Datengestützte Unternehmenssteuerung      \\

    \ac{HR}-Transformation      & Lernplattform, Talentmanagement           & Mitarbeiterentwicklung und Qualifizierung \\
    \hline
\end{xltabular}

Die Lernplattform ist zentrale Infrastruktur für künftige Weiterbildungs- und Qualifizierungsinitiativen.

\section{Projektpriorisierung im Portfolio}

Die Priorisierung erfolgt nach strategischer Relevanz, Nutzenbeiträgen und Umsetzbarkeit. Das Projekt „Lernplattform“ wurde wegen seines Beitrags zur Personalentwicklung und des wirtschaftlichen Potenzials (Reduktion administrativer Aufwände, Transparenz über Weiterbildung und Kompetenzen) priorisiert.

\begin{xltabular}{\textwidth}{@{}lcX@{}}
    \caption{Bewertung im Rahmen der Portfolio-Priorisierung} \label{tab:portfolio-priorisierung} \\
    \hline
    \textbf{Kriterium}       & \textbf{Bewertung} & \textbf{Begründung}                                                 \\
    \hline
    Strategischer Beitrag    & Hoch               & Unterstützt die Personalentwicklungsstrategie der dm-Gruppe         \\

    Wirtschaftlicher Nutzen  & Hoch               & Reduktion administrativer Aufwände im Schulungsmanagement           \\

    Organisatorischer Nutzen & Hoch               & Verbesserte Transparenz über Weiterbildungs- und Kompetenzdaten     \\

    Risiko                   & Mittel             & Organisationsweite Einführung und Integration bestehender Systeme   \\

    Umsetzungsaufwand        & Mittel             & Einführung eines konzernweiten Systems mit Change-Management-Anteil \\
    \hline
\end{xltabular}

Das Projekt wurde als strategische Initiative priorisiert und zur Umsetzung freigegeben.

\section{Portfolio-Governance}

Die Steuerung erfolgt in der Portfolio-Governance der dm-Gruppe. Das Steering Committee (\ac{HR}, \ac{IT}, Unternehmensleitung) trägt die Verantwortung, überwacht Fortschritt, Budget und Zielerreichung und trifft Anpassungsentscheidungen. Wichtige Entscheidungs- und Kontrollpunkte:

\begin{itemize}
    \item Freigabe des Business Case und der Projektinitialisierung
    \item Architektur- und Integrationsreview
    \item Pilotphase und Validierung der Plattform
    \item Entscheidung über den konzernweiten Rollout
\end{itemize}

So wird das Projekt im Portfolio bewertet und unterstützt die strategischen Ziele der dm-Gruppe.